\documentclass[twocolumn]{iopjournal}

\usepackage{amsmath,amssymb,amsfonts}

\fancyhfoffset[L]{0mm}
\setlength{\marginparwidth}{0pt}
\setlength{\marginparsep}{0pt}

% ==========================================
% ZENODO PREPRINT HEADERS (ACTIVE)
% ==========================================
\renewcommand{\articletype}[1]{%
  {\vspace*{-8mm}\noindent \Large \sf Preprint}\par
  \vspace*{6mm}\noindent{\scriptsize \sf{\bfseries \MakeUppercase{#1}}}}

\fancyhead[L]{}
\fancyhead[R]{}

\begin{document}

\articletype{Letter}

\title{Field Closure of General Relativity}

\author{Stephen Nelson}

\affil{Field Dynamics Research}

\ead{stephen@fielddynamics.org}

\begin{abstract}
The companion letter~\cite{Nelson2026b} derived the
exchange-symmetric partition $\mu(x) = x/(x{+}1)$ between the
two tetrahedral chambers of the closed frame, where
$x = g/a_0$ is the dimensionless acceleration.  That derivation
was algebraic: it determined the partition function but not its
spatial form.  This letter achieves field closure: it
presents the spatial structure of the complete 64-component
gravitational connection and closes the chain from geometry to
dynamics.
Evaluated radially, the partition produces
a continuously graded field whose inner throat is
Christoffel-dominated ($\mu \to 1$), whose outer disk is
torsion-dominated ($\mu \to 0$), and whose crossover ring at
$x = 1$ is the MOND transition scale $a_0$.  The monotonicity
of the partition produces a unique bulge-to-disk ordering for
any bound system within the law's domain.  The
24~torsion components organise into eight null directions
($2K = 8$ vertices of the closed frame, $D = 3$ rotation
planes each), whose chiral holonomy defines a rotation axis
and forces the field into a disk.  The radial acceleration relation
$g_{\mathrm{bar}} = g_{\mathrm{obs}}^2/(g_{\mathrm{obs}} +
a_0)$ follows directly from the partition.  Integration via the
AQUAL identity yields the unique scalar-sector energy
functional $F(y) = y - 2\sqrt{y} + 2\ln(1{+}\sqrt{y})$,
whose three terms correspond to the three structural levels of
the coupling polynomial.  The constitutive law determines the
action; the action determines the dynamics; no step in this
chain involves a free parameter.  General relativity's
torsion-free condition projects the 64-component connection
onto a 40-component subspace; the apparent separation of
baryonic and dark components is an artefact of this
incomplete projection.  Restoring the full connection closes
the field.
\end{abstract}

% ==================================================================

\section{The partition in space}

The companion letter~\cite{Nelson2026b} proved that the
exchange symmetry between TetA and TetB, combined with
rational bijectivity, uniquely determines the partition of
the gravitational field response between the two tetrahedral
chambers:
\begin{align}
  \text{TetA:} \quad & \mu(x) = \frac{x}{x + 1},
  \label{eq:mu} \\[4pt]
  \text{TetB:} \quad & 1 - \mu(x) = \frac{1}{x + 1},
  \label{eq:tetb}
\end{align}
where $x = g/a_0$ is the dimensionless acceleration.  The
derivation used no field equation and no fitting.  It required
only the parity relation between TetA and
TetB~\cite{Nelson2026a} together with the rational-bijection
classification on the projective line.  The present letter does
not re-derive this result.  It takes the function as given and
asks: what does this partition look like when it is evaluated
radially and revolved about a field axis?

The answer is Figure~\ref{fig1}.

The constitutive law $\mu(x) = x/(x{+}1)$ is a single
function.  But the tetrad equation
$g_{\mu\nu} = \eta_{ab}\,e^a{}_\mu\,e^b{}_\nu$ is quadratic
in the frame field~\cite{Nelson2026a}.  It has two roots in
the parity direction: TetA and TetB.  The constitutive law
therefore partitions into two complementary functions, one per
root:
\begin{center}
TetA receives $\mu(x)$. \quad TetB receives $1 - \mu(x)$.
\end{center}
At every radius, both roots are present; the two functions
specify how much of the total field response each root
carries.  They sum to unity everywhere.

Figure~\ref{fig1} renders each partition function as its own
surface.  The construction is direct: at each point in the
radial domain, $\mu(x)$ sets the height of the blue surface,
and $1 - \mu(x)$ sets the height of the red surface.  Both
surfaces share the same radial base.  The result is a
two-branch rendering over the radial domain --- one surface per
root of the quadratic equation --- with the vertical axis
measuring the field fraction carried by that root.  The two
surfaces cross where their fractions are equal, at $x = 1$
($\mu = 1/2$).  Figure~\ref{fig1} is therefore not a
heuristic sketch but the constitutive law itself rendered over
radius: once $\mu$ and $1 - \mu$ are plotted on the same
domain, the bulge-to-disk ordering follows directly.

A spiral galaxy has a bright central bulge surrounded by an
extended disk.  Why does Figure~\ref{fig1} reproduce that
structure?

Every gravitationally bound system has an acceleration
profile that is strongest at the centre and weakens outward.
This is not a model; it is Gauss's law applied to any finite
mass distribution.  The constitutive law
$\mu(x) = x/(x{+}1)$ is a monotonic function: higher
acceleration always maps to a larger TetA fraction, and lower
acceleration always maps to a larger TetB fraction.  The
composition of a monotonic partition with a monotonically
decreasing acceleration profile produces the same
bulge-to-disk radial ordering.  Where $x$ is large (the centre), TetA dominates
and the blue surface peaks.  Where $x$ is small (the rim),
TetB dominates and the red surface rises.  The result is a
peaked central concentration inside an extended disk --- a
bulge inside a disk --- for any gravitationally bound system,
not only for a particular galaxy model.  The resemblance
between Figure~\ref{fig1} and a spiral galaxy is therefore
not a visual coincidence but the direct rendering of the
constitutive law over radius.

\begin{figure}[t]
\centering
\includegraphics[width=\columnwidth]{fig1_tetrad_partition}
\caption{The two branches of the quadratic tetrad equation,
projected through the constitutive law $\mu(x) = x/(x{+}1)$
and displayed as surfaces over the radial domain.  Blue: the
TetA branch $\mu$ (the Christoffel sector retained by general
relativity).  Red: the TetB branch $1 - \mu$ (the torsion
sector removed by the zero-torsion condition).  At every
radius, blue~$+$~red~$= 1$.  The white ring marks the
crossover at $x = 1$, where both roots contribute equally
($\mu = 1/2$).  Their overlap at centre corresponds to the
galactic bulge; the extended red surface represents the field
fraction that sustains flat rotation
at the rim.  Eight spiral arms alternate between TetA (blue)
and TetB (amber), corresponding to the $2K = 8$ null
directions of the closed frame.}
\label{fig1}
\end{figure}

The vertical axis in Figure~\ref{fig1} is the field fraction,
not a physical height.  The blue spike at centre is TetA
carrying nearly the entire field response.  The red plateau at
large radius is TetB carrying nearly all of it.  Between them,
the white ring marks the handover.  The textbook decomposition
of the full connection into Christoffel plus contorsion
corresponds, at the level of field magnitudes, to the
partition into $\mu$ and
$1 - \mu$~\cite{Nelson2026b}.  Figure~\ref{fig1} is that
decomposition rendered spatially.

The white ring sits at $x = 1$, where $\mu = 1/2$.  This is
the locus where both chambers contribute equally to the
field response.  Inside the ring, the Christoffel sector
dominates: general relativity is accurate and Keplerian
decline holds.  Outside the ring, the torsion sector
dominates: the field response that general relativity
discards becomes the majority.  This crossover is the MOND
transition scale $a_0$~\cite{Milgrom1983}.  It is not a
fitted parameter.  It is the geometric midpoint of the
exchange-symmetric partition, the point where the odds ratio
$\Omega = \mu/(1-\mu)$ equals unity.

The companion letter~\cite{Nelson2026b} presented the three
acceleration regimes in tabular form: strong field
($x \gg 1$, $\mu \approx 1$), transition ($x = 1$,
$\mu = 1/2$), and weak field ($x \ll 1$, $\mu \approx 0$).
Figure~\ref{fig1} gives those three regimes spatial identity:
the throat, the ring, and the disk.

The radial acceleration
relation~\cite{McGaugh2016,Li2018} is reproduced by this
partition and rendered here as geometry rather than as a
scatter plot.  The blue
spike represents the baryonic contribution.  The red surface
represents the field fraction whose effects standard
cosmological analysis typically models with an invisible mass
component.  They sum to one everywhere.

% ==================================================================

\section{The radial transition}

The full affine connection on the tetrad frame has $K^3 = 64$
independent
components~\cite{Nelson2026b,Hehl1976}.  The full
64-component connection is \textit{field-closed}: both roots of the tetrad
equation contribute, the partition $\mu + (1{-}\mu) = 1$ holds
at every radius, and the field response is
self-consistent without external mass.  This is distinct from
Einstein's notion of a ``complete field theory'' (in which the
equations of motion follow from the field equations without
independent postulates) and from geodesic completeness (the
extendibility of geodesics to infinite affine parameter).
Field closure refers specifically to the connection: whether
all $K^3$ components of the affine connection are active, or
whether a subset has been set to zero by constraint.
Einstein--Cartan theory~\cite{Hehl1976} restores torsion but
couples it algebraically to spin, so that it vanishes in
vacuum.  The present framework restores torsion with an
acceleration-dependent partition: the 24~antisymmetric
components carry the fraction $1 - \mu(x)$ of the field
response, which is nonzero wherever $x$ is finite.

General relativity's
torsion-free condition projects this onto the 40~symmetric
components (Christoffel), discarding the 24~antisymmetric
components (torsion).  With only 40 of 64 components, the
connection is not field-closed.  In Figure~\ref{fig1}, this projection
removes the red surface.  What remains is the blue surface
alone: a central concentration that peaks at high acceleration
and declines to zero at the rim.  Without the red surface, the
field cannot sustain flat rotation at large radius, because
the TetA fraction $\mu = x/(x{+}1)$ vanishes where the
rotation curves are flat.  The 24~torsion components carry the
field response of TetB, the second root, which rises precisely
where TetA declines.  Their overlap at the core corresponds
to the galactic bulge.  The crossover ring at $x = 1$ marks
the transition between the two
regimes~\cite{FamaeyMcGaugh2012}.

The rate at which the partition shifts between the two roots
is characterised by a logistic identity that $\mu$
satisfies~\cite{Nelson2026b}:
\begin{equation}
  x\,\frac{d\mu}{dx} = \mu(1 - \mu).
  \label{eq:logistic}
\end{equation}
Under the substitution $t = \ln x$, this becomes
$d\mu/dt = \mu(1-\mu)$: the canonical two-sector transition
law.  The product $\mu(1-\mu)$ on the right-hand side measures
the simultaneous presence of both tetrahedral chambers at a
given radius.  It is the continuous analogue of the octahedral
interaction core, the $s = 0$ state of the coupling
polynomial~\cite{Nelson2026b}.  The equation states that the
partition can shift only where both chambers contribute, which
is to say, only where the interaction core is active.  At the
deep centre, one chamber dominates and the partition is frozen.
At the rim, the same.  The transition peaks at the MOND ring
because that is where the core is fully realised:
$\mu(1 - \mu) = 1/4$.

The supplementary note to this letter establishes the
geometric origin of this constraint: every transition between
a TetA and a TetB null direction crosses a face of the
octahedral core, generating the chiral rotation that forces
the field into a disk.  The logistic equation is the
continuous, radial expression of that discrete crossing
requirement.

The apparent separation of baryonic and dark components in
Figure~\ref{fig1} is a diagnostic decomposition of one field,
not a claim of two independent physical substances.  It is the
projection of a single 64-component field onto a 40-component
subspace.  The observed gravitational discrepancy at galactic
scales is not in dispute: both the cold dark matter framework
and the present partition reproduce the radial acceleration
relation and the function $\mu(x) = x/(x{+}1)$.  They differ
on what the discrepancy represents.  The cold dark matter
framework interprets it as an invisible mass component,
modelled with approximately spherical
halos~\cite{NFW1996}.  The closed tetrad frame identifies it as the field
response of the second tetrahedral chamber, carried by the
24~torsion components.  This response is not mass.  It is
topology.  Its spatial distribution is not fitted but
determined: the logistic equation~(\ref{eq:logistic}) governs
the radial profile, and the result is the extended disk
visible as the red surface in Figure~\ref{fig1}.

% ==================================================================

\section{The eight-channel structure}

Figure~\ref{fig1} displays eight spiral
arms, four coloured blue (TetA) and four amber (TetB),
alternating around the disk.  This section identifies what they
are and why they force the field into a disk rather than a
sphere.

The companion letter~\cite{Nelson2026b} proved (Section~5,
Theorem) that the 24~torsion components admit the
factorisation
\begin{equation}
  24 \;=\; \underbrace{8}_{\text{null directions}}
        \;\times\; \underbrace{3}_{\text{rotation planes}}
  \;=\; 2K \times D
  \label{eq:24}
\end{equation}
exclusively at $K = 4$.  In the companion letter's polynomial
language, $2K$ is the coupling range
$\Delta f = f(K,{+}1) - f(K,{-}1)$.  The algebraic identity
was established there.  The supplementary note to the present
letter identifies the two factors geometrically.  The
factorisation $8 \times 3$ is not one of several ways to
decompose 24; it is the unique decomposition in which one
factor is the vertex count of the closed frame and the other
is the spatial dimension.

The factor $2K = 8$ is the count of null directions of the
closed frame: the eight vertices of the stellated
octahedron~\cite{Nelson2026a}, equivalently the eight vertices
of the cube dual to the octahedral core.  Each face of the
octahedral core carries an attached tetrahedron whose apex is
one of these eight vertices.  The base face, shared with the
octahedron, is the propagation face of the corresponding channel.  The three lateral faces are the three transverse
rotation planes.  The eight spiral arms in
Figure~\ref{fig1} are these eight null directions
projected onto the disk.  Four belong to TetA.  Four belong to
TetB.  They alternate because adjacent faces of the octahedral
core always carry opposite chirality.

The supplementary note proves that parallel transport through
all eight null directions produces trivial holonomy when the
TetA/TetB chirality distinction is ignored, but nontrivial
holonomy when it is respected.  All twelve Hamiltonian cycles
on the cube graph were computed.  In the uniform case (all
rotations $+\theta$), all twelve yield the identity.  In the
chiral case (rotation sign determined by the source vertex's
sub-frame), all twelve yield nontrivial rotations: six at
$135^{\circ}$ and six at $155.8^{\circ}$.  The nontrivial
holonomy defines a rotation axis, and because the eight
channels all carry transport through the common field origin
with a consistent rotational sense, it organises the field
into a rotationally coherent disk perpendicular to that axis.
Because the transport accumulates into shared planar
circulation rather than cancelling into spherical closure, the
partition forms a disk and not a sphere.  The
axis is not imposed by hand.  It emerges from the chirality
of the closed frame.  The bridge from local frame structure to
macroscopic geometry is the constitutive law itself: $\mu(x)$
is a function of the acceleration, a base-manifold quantity,
and it distributes the field response between the two roots at
every spacetime point.  Because the torsion sector carries an
intrinsic axial preference at every point, the macroscopic
field built from it inherits that preference globally.

The kinematic consequence follows directly from the partition.
At dimensionless acceleration $x = g_{\mathrm{obs}}/a_0$, the
baryonic acceleration, carried by the Christoffel sector, is
$g_{\mathrm{bar}} = \mu\,g_{\mathrm{obs}}$, giving
\begin{equation}
  g_{\mathrm{bar}} =
  \frac{g_{\mathrm{obs}}^2}{g_{\mathrm{obs}} + a_0}.
  \label{eq:rar}
\end{equation}
At high acceleration ($x \gg 1$):
$g_{\mathrm{bar}} \approx g_{\mathrm{obs}}$, Keplerian
decline, general relativity is accurate.  At low acceleration
($x \ll 1$):
$g_{\mathrm{bar}} \approx g_{\mathrm{obs}}^2/a_0$, flat
rotation curves.  The radial acceleration relation established
empirically by McGaugh, Lelli \&
Schombert~\cite{McGaugh2016}, with 0.057~dex scatter across
175~individual galaxies~\cite{Li2018} and no credible
variation in $a_0$, is reproduced by Eq.~(\ref{eq:rar}) with
no free parameters.  The full 179-galaxy rotation curve
analysis is presented in a companion
paper~\cite{Nelson2026c}.

% ==================================================================

\section{What the partition resolves}

The spatial partition developed in Sections~1--3 makes three
observational predictions that distinguish it from the cold
dark matter framework.

\textit{Disk geometry.}  Disk galaxies typically consist of a
stellar disk, a gas disk, and a central bulge.  The cold dark matter
framework adds a spherical halo to account for the
gravitational discrepancy at large
radius~\cite{NFW1996}.  The closed tetrad frame predicts that
the additional field response is also disk-like, not spherical,
because the chiral holonomy of the eight-channel structure
defines an axis and same-sense transport organises the field
into a planar disk perpendicular to it (supplementary note).
No halo geometry is fitted.  Because the constitutive law
modulates the effective torsion strength radially, the chiral
holonomy is suppressed where $\mu \approx 1$ (the centre) and
active where $\mu \approx 0$ (the rim).  The partition
therefore produces not only a disk but a spherically symmetric
bulge transitioning smoothly into a flat disk, with the
morphological transition at the MOND ring $x = 1$.

\textit{Bounded central density.}  The TetB field fraction
$1/(x{+}1)$ is bounded above by unity for all $x$.  The
torsion-sector field acceleration never exceeds $a_0$.  The NFW
profile diverges as $r \to 0$~\cite{NFW1996}.  The partition's
boundedness removes the central divergence responsible for the
cusp-core tension, without baryonic
feedback or self-interacting dark matter.  The removal is
algebraic: $1/(x{+}1) < 1$ is a property of the constitutive
law, not of any galaxy model.

\textit{No new particles.}  The gravitational discrepancy is
not mass.  It is the field response of the second tetrahedral
chamber, carried by the 24~torsion components that the
zero-torsion condition removes from the connection.  The
missing component is not in the matter sector.  It is in the
connection.

General relativity succeeds in the solar system because the
solar system sits deep inside a galaxy where the gravitational
acceleration exceeds the transition scale $a_0$ by a factor of
order $10^8$.  At those accelerations the partition fraction
$\mu$ is so close to unity that the second root contributes
less than one part in a million.  The torsion-free
approximation loses nothing measurable, and general relativity
was confirmed on every local test it was given.  But a galaxy
spans the full range of the partition: from the core where one
root carries nearly everything, through the crossover ring
where both roots contribute equally, to the rim where the
second root carries the majority.  The function
$\mu(x) = x/(x{+}1)$ connects these regimes smoothly.
Milgrom~\cite{Milgrom1983} identified it empirically.  Famaey
\& McGaugh~\cite{FamaeyMcGaugh2012} selected it by fitting
galactic data.  The constitutive law that appears
phenomenologically in the AQUAL formulation of modified
gravity is here recovered as the geometric partition of the
two-chamber tetrad frame.  The present series derives it: the
tetrad equation is degree two, exchange symmetry between its
roots admits exactly one rational partition, and that
partition is $\mu$.

The constitutive law also determines the scalar-sector
energy functional.  The variational relationship between a
constitutive law and its energy functional is standard
Lagrangian mechanics: the response function is the derivative
of the potential.  Applied to gravity, this is the identity
$\mu(\sqrt{y}) = dF/dy$ used by Bekenstein \&
Milgrom~\cite{BekensteinMilgrom1984} in the AQUAL
formulation, where
$y = |\nabla\Phi|^2/a_0^2 = x^2$ is the squared
dimensionless acceleration.  Integration of
$\mu(x) = x/(x{+}1)$ yields the unique kinetic functional
\begin{equation}
  F(y) = y - 2\sqrt{y} \;+\; 2\ln\!\bigl(1 + \sqrt{y}\bigr),
  \label{eq:F}
\end{equation}
In the standard AQUAL literature, $F$ is a free function
selected to fit observations.  Here it is fixed by the
constitutive law, which is itself fixed by the frame geometry.
The three terms of Eq.~(\ref{eq:F}) correspond to the three
structural levels of
the coupling polynomial~\cite{Nelson2026b}: the quadratic $y$
to the $K^2$ face-coupling level, the $-2\sqrt{y}$ to the
state-dependent $sK$ propagation level, and the logarithm to
the $K^0$ field origin.  At high acceleration ($y \gg 1$),
$F \to y$ and the action reduces to general relativity.  At
low acceleration ($y \ll 1$), $F \to \tfrac{2}{3}\,y^{3/2}$,
recovering the deep-field regime and the baryonic Tully-Fisher
relation.  The constitutive law determines the action.  The
action determines the dynamics.  No step in this chain
involves a free parameter.

To summarise the downstream structure: the partition
$\mu = x/(x{+}1)$ is selected uniquely by exchange symmetry
between the two tetrad roots~\cite{Nelson2026b}.  The same
function partitions the 64-component connection into the
40~Christoffel and 24~torsion sectors.  Its complement
$1 - \mu$ is exactly the field fraction that, in the
torsion-free projection, must be attributed to invisible mass.
It satisfies the logistic transition equation
$d\mu/d\!\ln x = \mu(1{-}\mu)$.  Its integral yields the
unique AQUAL energy functional $F(y)$, whose three terms map
to the three structural levels of the coupling polynomial.
That functional recovers general relativity at high
acceleration and the baryonic Tully-Fisher relation at low
acceleration.  These are not separate results.  They are the
same algebraic object encountered at successive levels of the
theory: partition, connection, transition, action, kinematics.
These closures are not independent.  Rejection of any one
therefore requires an alternative account of the whole chain.

An alternative relativistic approach to the MOND phenomenology
has been developed by Skordis \&
Z\l{}o\'{s}nik~\cite{SkordisZlosnik2021}, who introduce
additional vector and scalar fields to obtain a theory
compatible with the cosmic microwave background power spectrum.
The present framework adds no fields.  It restores the
24~connection components that general relativity's zero-torsion
condition removes, recovering the complete geometry of the
tetrad frame.  The result is field closure: the gravitational
connection carries all $K^3 = 64$ components, the two
tetrahedral chambers contribute their full partition at every
radius, and the field response balances without invisible mass.

The torsion-free specialisation adopted in 1915 discards the
24~antisymmetric components of the affine connection.
Restoring them closes the field.  The field response they
carry is precisely what, in their absence, must be attributed
to invisible mass.

% ==================================================================

\section{Data and reproducibility}

The chiral holonomy computation is provided as a supplementary
note (\texttt{supplementary\_chiral\_holonomy.pdf}).  A
self-contained verification script
(\texttt{holonomy\_verify.py}, approximately forty lines of
Python using NumPy) reproduces the holonomy of all twelve
Hamiltonian cycles on the cube graph for both uniform and
chiral transport.  An interactive rendering of
Figure~\ref{fig1} is provided as a
supplementary HTML file permitting rotation, zoom, and
background toggling.  A second interactive file illustrates
the eight-channel vortex structure of Section~3.  To complete
the field-closure claim, a second supplementary note
(\texttt{supplementary\_covariant\_completion.pdf}) derives the
full four-term covariant action from the constitutive law,
confirming that the closed connection yields a complete
relativistic field theory with no free functions.

% ==================================================================

\begin{thebibliography}{99}

\bibitem{Nelson2026a}
S.~Nelson,
\textit{The Coordinate-Space Geometry of the Tetrad Frame},
Class.\ Quantum Grav.\ (2026), submitted (CQG-115428).

\bibitem{Nelson2026b}
S.~Nelson,
\textit{The Other Half of Gravity: Torsion and the
Two-Chamber Structure of the Tetrad Frame},
Class.\ Quantum Grav.\ (2026), submitted (CQG-115447).

\bibitem{Nelson2026c}
S.~Nelson,
\textit{Gravity Field Dynamics: Decoding mass from observed
rotation curves},
Mon.\ Not.\ R.\ Astron.\ Soc.\ (2026), submitted.

\bibitem{McGaugh2016}
S.~S.~McGaugh, F.~Lelli and J.~M.~Schombert,
\textit{Radial Acceleration Relation in Rotationally
Supported Galaxies},
Phys.\ Rev.\ Lett.\ \textbf{117} (2016) 201101.

\bibitem{Li2018}
P.~Li, F.~Lelli, S.~McGaugh and J.~Schombert,
\textit{Fitting the radial acceleration relation to
individual SPARC galaxies},
Astron.\ Astrophys.\ \textbf{615} (2018) A3.

\bibitem{Milgrom1983}
M.~Milgrom,
\textit{A modification of the Newtonian dynamics as a
possible alternative to the hidden mass hypothesis},
Astrophys.\ J.\ \textbf{270} (1983) 365.

\bibitem{FamaeyMcGaugh2012}
B.~Famaey and S.~S.~McGaugh,
\textit{Modified Newtonian Dynamics (MOND): Observational
Phenomenology and Relativistic Extensions},
Living Rev.\ Relativity \textbf{15} (2012) 10.

\bibitem{BekensteinMilgrom1984}
J.~Bekenstein and M.~Milgrom,
\textit{Does the missing mass problem signal the breakdown
of Newtonian gravity?},
Astrophys.\ J.\ \textbf{286} (1984) 7.

\bibitem{Hehl1976}
F.~W.~Hehl, P.~von~der~Heyde, G.~D.~Kerlick and
J.~M.~Nester,
\textit{General relativity with spin and torsion:
Foundations and prospects},
Rev.\ Mod.\ Phys.\ \textbf{48} (1976) 393.

\bibitem{NFW1996}
J.~F.~Navarro, C.~S.~Frenk and S.~D.~M.~White,
\textit{The structure of cold dark matter halos},
Astrophys.\ J.\ \textbf{462} (1996) 563.

\bibitem{SkordisZlosnik2021}
C.~Skordis and T.~Z\l{}o\'{s}nik,
\textit{New Relativistic Theory for Modified Newtonian
Dynamics},
Phys.\ Rev.\ Lett.\ \textbf{127} (2021) 161302.

\end{thebibliography}

\end{document}
